% =========================================================
\section{实验结果}
\label{sec:results}
% =========================================================

\subsection{主要检索结果}

表~\ref{tab:retrieval_results} 给出了测试集上的检索性能。
\REV{TEF-RAG 在 Recall@5、Complete@5 和 FlowComplete@5 上取得最高值，而 Fine-tuned BGE Reranker 在 nDCG@5 上取得最高值。与 Fine-tuned BGE 相比，TEF-RAG 将 Recall@5 从 0.6650 提升至 0.7228，将 Complete@5 从 0.2854 提升至 0.3188，并将 FlowComplete@5 从 0.2833 提升至 0.3146；另一方面，Fine-tuned BGE 的 nDCG@5 更高（0.6501 对 0.6231）。这一现象与单条证据排序质量和完整证据集合恢复能力之间的区别一致。}

\begin{table}[t]
\caption{测试集上的主要检索结果（Top-5）。}
\label{tab:retrieval_results}
\centering
\small
\setlength{\tabcolsep}{3pt}
\resizebox{\columnwidth}{!}{%
\begin{tabular}{lrrrr}
\toprule
方法 & Recall & nDCG & Complete & FlowComplete \\
\midrule
BM25 & 0.3791 & 0.3915 & 0.0563 & 0.0563 \\
BGE Reranker & 0.2916 & 0.3070 & 0.0417 & 0.0417 \\
\rowcolor{yellow}
Fine-tuned BGE & 0.6650 & \textbf{0.6501} & 0.2854 & 0.2833 \\
Temporal-BM25 & 0.4591 & 0.4955 & 0.0792 & 0.0792 \\
TA-RAG & 0.2769 & 0.2907 & 0.0375 & 0.0375 \\
\textbf{TEF-RAG} & \textbf{0.7228} & \cellcolor{yellow}0.6231 & \textbf{0.3188} & \textbf{0.3146} \\
\bottomrule
\end{tabular}%
}
\end{table}

\subsection{验证集消融实验}
\label{sec:ablation_results}

表~\ref{tab:ablation_results} 给出了验证集上的四组消融结果。完整 TEF-RAG 的 FlowComplete@5 为 0.4042。移除学习式证据对提议后，该指标仅下降至 0.4021。在送入关系判断的证据对数量保持不变的条件下，这一较小变化表明，证据对提议器主要决定哪些证据对接受关系分析，而不是直接决定最终证据集合。相比之下，移除由关系图导出的特征后，FlowComplete@5 降至 0.3667；这一下降与图结构特征对集合排序存在贡献的预期一致。

训练目标与排序器本身的影响更大。移除难负样本挖掘后，FlowComplete@5 从 0.4042 降至 0.2833；移除非线性集合排序器、改用前序贪心证据集合选择器后，该指标降至 0.3125。在这一固定验证设置下，移除难负样本与移除非线性排序器所造成的 FlowComplete@5 下降均大于移除证据对提议的影响。值得注意的是，相比完整模型，移除非线性排序器后 nDCG@5 略有上升（0.6965 对 0.6885），但 Complete@5 与 FlowComplete@5 明显下降。nDCG@5 上升而 Complete@5 与 FlowComplete@5 下降也说明，传统排序质量与集合级过程完整性并不一定沿同一方向变化。

\begin{table}[t]
\caption{TEF-RAG 在验证集上的消融结果。}
\label{tab:ablation_results}
\centering
\small
\setlength{\tabcolsep}{3pt}
\resizebox{\columnwidth}{!}{%
\begin{tabular}{lrrrr}
\toprule
变体 & Recall & nDCG & Complete & FlowComplete \\
\midrule
完整 TEF-RAG & 0.7567 & 0.6885 & 0.4063 & 0.4042 \\
移除 Pair Proposal & 0.7494 & 0.6832 & 0.4063 & 0.4021 \\
移除 Graph Features & 0.7184 & 0.6601 & 0.3750 & 0.3667 \\
移除 Hard Negatives & 0.6748 & 0.6042 & 0.2875 & 0.2833 \\
移除 Nonlinear Ranker & 0.7070 & 0.6965 & 0.3188 & 0.3125 \\
\bottomrule
\end{tabular}%
}
\end{table}

\subsection{下游结构化生成结果}
\label{sec:generation_results}

表~\ref{tab:generation_results} \REV{给出了四个互补的生成指标，用于比较不同检索证据所产生的结构化输出。在生成器、提示词、输出格式和解码设置均相同的情况下，TEF-RAG 在四个展示指标上均取得最高值：Field Macro F1 为 0.5345，Action F1 为 0.2035，Citation F1 为 0.1585，Evidence Support Recall 为 0.2253。对 BGE 进行任务内微调后，其结果相比预训练版本明显提高，但 TEF-RAG 在这四项指标上仍保持更高值。结果表明，TEF-RAG 所检索的证据不仅反映到字段级与动作级输出质量上，也反映到 claim--evidence 引用质量，以及参考 claim 被可接受支撑证据覆盖的程度上。}

\begin{table}[t]
\caption{生成测试集上的主要结构化生成指标。}
\label{tab:generation_results}
\centering
\small
\setlength{\tabcolsep}{3pt}
\resizebox{\columnwidth}{!}{%
\begin{tabular}{lrrrr}
\toprule
方法 & Field F1 & Action F1 & Citation F1 & \cellcolor{yellow}Evidence Support Recall \\
\midrule
BM25 & 0.4715 & 0.1372 & 0.0813 & \cellcolor{yellow}0.1258 \\
BGE Reranker & 0.4129 & 0.1198 & 0.0405 & \cellcolor{yellow}0.0614 \\
\rowcolor{yellow}
Fine-tuned BGE & 0.4818 & 0.1910 & 0.0954 & 0.1328 \\
Temporal-BM25 & 0.5218 & 0.1637 & 0.1249 & \cellcolor{yellow}0.1871 \\
TA-RAG & 0.4059 & 0.1115 & 0.0379 & \cellcolor{yellow}0.0570 \\
\textbf{TEF-RAG} & \textbf{0.5345} & \textbf{0.2035} & \textbf{0.1585} & \cellcolor{yellow}\textbf{0.2253} \\
\bottomrule
\end{tabular}%
}
\end{table}

% =========================================================
\section{局限性与有效性威胁}
\label{sec:limitations}
% =========================================================

本文基准是一个以公开资料为依据、由 AI 辅助构建的合成基准，而不是真实电站工单集合，也不是领域专家人工标注的数据集。
因此，其主要价值在于对时序证据检索机制进行受控比较；实验结果不应被直接外推为真实电站的故障频率或现场机器人的安全性结论。
尽管生成评估覆盖了字段、动作、依赖关系与引用，但并未模拟真实机器人执行、环境反馈或安全联锁。

% =========================================================
\section{结论}
\label{sec:conclusion}
% =========================================================

本文提出 TEF-RAG，将动态运维 RAG 的检索目标从独立文档相关性扩展为在查询时刻可见、关系一致且过程完整的时序证据流。
\REV{在测试集上，六种对比方法中，TEF-RAG 在 Recall@5、Complete@5 与 FlowComplete@5 上取得最高值，而 Fine-tuned BGE 在 nDCG@5 上取得最高值。在固定下游生成器的结构化生成测试中，TEF-RAG 在 Field Macro F1、Action F1、Citation F1 与 Evidence Support Recall 上均取得最高值。}
未来工作应进一步整合显式的规划与执行约束。
